% =====================================================================
\section{Predicted Problems}\label{ch:predicted}
% =====================================================================

Part A is a full predicted quiz in exactly the Quiz~2 format and difficulty,
now including Slides~7. Part B gives extra drills on each pattern. Solutions
are in Chapter~\ref{ch:predsol}.

\subsection{Part A: predicted quiz (55 points)}

\begin{questions}
\item Let $C$ be the curve of intersection of $x^2+y^2+z^2=100$ and $4x-3z=0$.
      \begin{enumerate}[label=(\alph*)]
        \item \textbf{[6 points]} Find $\vr(t)$ parametrizing $C$, oriented counterclockwise
              when viewed from above, starting at $P(6,0,8)$.
        \item \textbf{[4 points]} Find a vector equation of the tangent line to $C$ at
              $R\left(3,\,5\sqrt3,\,4\right)$.
        \item \textbf{[5 points]} Find the exact arc length of $C$ from $P(6,0,8)$ to
              $S(-6,0,-8)$, following the orientation.
        \item \textbf{[7 points]} Find the curvature of $C$ at $P$.
      \end{enumerate}\label{pq:1}
\item Let $f(x,y)=3xy^2+\ln(x^2+y^2)-\dfrac{2y}{x}$, $P(1,1)$, $Q(4,-3)$.
      \begin{enumerate}[label=(\alph*)]
        \item \textbf{[7 points]} Find the directional derivative of $f$ at $P$ toward $Q$.
        \item \textbf{[4 points]} Find the unit vector in the direction of fastest increase at $P$.
        \item \textbf{[2 points]} Find the maximum rate of increase at $P$.
        \item \textbf{[2 points]} Find the unit vectors along which the rate of change at $P$ is $0$.
      \end{enumerate}\label{pq:2}
\item Let $\vF(x,y,z)=\ip{y\cos(xy)+2xz,\;x\cos(xy)+e^z,\;x^2+ye^z}$.
      \begin{enumerate}[label=(\alph*)]
        \item \textbf{[4 points]} Show that $\vF$ is conservative.
        \item \textbf{[5 points]} Find the potential function $f$ with $f(0,0,0)=2$.
        \item \textbf{[4 points]} Find the maximum rate of increase of $f$ at $(1,0,0)$ and
              the directional derivative of $f$ there toward $(1,2,2)$.
      \end{enumerate}\label{pq:3}
\end{questions}

\textbf{Problems 4--8: answer any; your best TWO count (5 points each).}

\begin{questions}[resume]
\item Start at $(1,0,1)$ and move $\sqrt3$ units along
      $\vr(t)=\ip{e^t\cos t,\;e^t\sin t,\;e^t}$ in the direction of increasing $t$.
      Find your new position.\label{pp:arc}
\item Find all $b$ such that the curve $y=b\sin x$ has curvature $2$ at $x=\frac\pi2$.
      Use a vector-valued parametrization.\label{pq:5}
\item Find the directions (unit vectors) in which the directional derivative of
      $f(x,y)=x^2+2xy$ at $(1,1)$ equals $2$.\label{pq:6}
\item Find all points at which the direction of fastest increase of
      $f(x,y)=x^2+2y^2-4x+4y$ is $\vi-\vj$.\label{pq:7}
\item Find constants $a,b,c$ such that
      $\vF=\ip{2xy+az^3,\;bx^2+3z,\;6xz^2+cy}$ is conservative, and find a potential.\label{pq:8}
\end{questions}

\subsection{Part B: extra drills}

\begin{questions}[resume]
\item Let $C$ be the intersection of the cylinder $4x^2+y^2=4$ with the plane $z=2$.
      (a) Parametrize $C$ counterclockwise from above, starting at $(1,0,2)$.
      (b) Find the maximum and minimum curvature of $C$ and where they occur.
      (c) Set up (do not evaluate) the integral for the length of $C$.\label{pd:1}
\item Start at $(2,0,0)$ and move $6\pi$ units along $\vr(t)=\ip{2\cos t,2\sin t,\sqrt5\,t}$
      in the direction of increasing $t$. Where do you end up?\label{pd:2}
\item At which point does the curve $y=e^x$ have maximum curvature? What is the
      maximum curvature?\label{pd:3}
\item Let $f(x,y,z)=xe^{yz}+\ln(x+y+z)$. Find the directional derivative at $(1,0,0)$
      toward $(3,2,1)$ and the maximum rate of change there.\label{pd:4}
\item Show that $\vF=\ip{\dfrac{2x}{1+x^2}+y^2,\;2xy+\dfrac1y}$ ($y>0$) is conservative and find a potential.\label{pd:5}
\item Show that $\vF=\ip{yz,\;xz,\;xy+x}$ is not conservative.\label{pd:6}
\item Let $f(x,y,z)=x^2y+yz^3$, $P(1,2,-1)$, $Q(2,4,1)$. (a) Find the gradient field
      $\vF=\grad f$ and verify it passes the component test. (b) Find the direction and
      value of the maximum rate of change of $f$ at $P$. (c) Find $\Du f(P)$ toward $Q$.\label{pd:7}
\item Describe the vector field $\vF(x,y)=\ip{y,x}$. Is it a gradient field?\label{pd:8}
\end{questions}
